\usepackage{amsmath, amsfonts, amssymb, amsthm, dsfont, mathrsfs}
\usepackage{adjustbox, graphicx, enumitem, fancyhdr, titlesec, etoolbox, environ, tabularx}
\usepackage[svgnames, dvipsnames]{xcolor}
\usepackage[left=1.25in, right=1.25in, top=1in, bottom=1in]{geometry}
\usepackage[inline]{asymptote}
\begin{asydef}
import olympiad;
defaultpen(fontsize(10pt));
usepackage("amsmath");
usepackage("amsfonts");
usepackage("amssymb");
\end{asydef}
\usepackage{imakeidx}
\usepackage{hyperref}
\usepackage[nameinlink]{cleveref}
\usepackage[framemethod=Tikz]{mdframed}
\mdfsetup{nobreak=false}

\usepackage[backend=biber, style=alphabetic, sorting=nyt]{biblatex}
\addbibresource{references.bib}

\usepackage{mark-macros}

\usepackage{tikz-cd}


%----------------------------------------------------%
%                 Formatting Setup
%----------------------------------------------------%
\usepackage[tt=false]{libertine}
\usepackage[scale=0.9]{noto-sans}

\fancyhf{}
\cfoot{-- \thepage~--}

\fancypagestyle{plain}{ %
  \fancyhf{} % remove everything
  \renewcommand{\headrulewidth}{0pt} % remove lines as well
  \cfoot{-- \thepage~--}
}

	
%heading setup

\newcommand{\todo}[1]{{\color{blue}[TODO: #1]}}
\setlength{\parskip}{1ex plus 0.5ex minus 0.2ex}
\setlength\arraycolsep{2pt}

\makeindex[name=notations, title=Notations, intoc]
\newcommand{\notation}[2]{\index[notations]{#1 {\small(#2)}}}

\titleformat{\chapter}[hang]{\Large\bfseries\sffamily}
{\scalebox{2.5}{\textcolor{BrickRed}{\thechapter}}\hspace{0.8cm}}{0pt}{\huge\bfseries\sffamily\boldmath}

\titleformat{\section}[hang]{\Large\bfseries\sffamily}{\textcolor{BrickRed}{$\boldsymbol{\S}$\thesection}\hspace{0.6cm}}{0pt}{\Large\bfseries\sffamily\boldmath}

\titleformat{\subsection}[hang]{\large\bfseries\sffamily}{\textcolor{BrickRed}{$\boldsymbol{\S}$\thesubsection}\hspace{0.6cm}}{0pt}{\large\bfseries\sffamily\boldmath}

\titleformat{\paragraph}[runin]{\bfseries\sffamily}{}{0pt}{$\blacktriangleright$ \bfseries\sffamily\boldmath}


\renewcommand{\textbf}[1]{{\bfseries\boldmath #1}}

\newenvironment{alphenum}{\begin{enumerate}[label=(\alph*)]}
{\end{enumerate}}

\pagestyle{fancy}
\fancyhead[LE, RO]{\thepage}
\fancyhead[RE]{\nouppercase{\rightmark}}
\fancyhead[LO]{\nouppercase{\leftmark}}
\fancyfoot[LE, LO, RE, RO, CE, CO]{}
\cfoot{-- \thepage~--}
\setlength\headheight{15pt}

\hypersetup{
	 colorlinks,
  	 linkcolor={MediumBlue!80!Black},
   	 citecolor={green!50!black},
  	 urlcolor={MediumBlue!80!Black}
}


%----------------------------------------------------%
%                indexing
%----------------------------------------------------%
\usepackage{ifthen}
\makeindex[intoc]
\newcommand{\vocab}[2][]{%
	\textcolor{BrickRed}{\textbf{\boldmath #2}}%
	\index{\ifstrempty{#1}{#2}{#1}}%
}


%----------------------------------------------------%
%                mdframed setup
%----------------------------------------------------%
\expandafter\let\expandafter\oldproof\csname\string\proof\endcsname
\let\oldendproof\endproof
\renewenvironment{proof}[1][\proofname]{%
\oldproof[\sffamily\itshape\color{MediumBlue!80!Black}$\blacktriangleright$ #1\color{black}]%
}{\oldendproof}

\usepackage[svgnames, dvipsnames]{xcolor}

\usepackage{twemojis}
\usepackage[skins, breakable]{tcolorbox}
\usepackage{keytheorems}
\tcbuselibrary{theorems}

\tcbset{enhanced,
	coltitle=black, 
	colback=white,
	boxsep=1.75mm, 
	sharp corners,
	toprule=0.5pt,
	leftrule=0pt,
	titlerule=0pt,
	bottomrule=0.5pt,
	rightrule=0.5pt,
	segmentation style={solid},
	fonttitle = \bfseries\sffamily,
	breakable=true,
	parbox=false
}
\newkeytheoremstyle{thmbox}
  {
    tcolorbox={
		colbacktitle=MediumBlue!15,
		colframe=MediumBlue!40,
		borderline west={1mm}{0pt}{MediumBlue!60!white},
		before title=\large\texttwemoji{pencil}\normalsize\hspace{2mm},
    },
    notebraces={(}{)},
    headpunct={.},
    noteseparator={ },
    headfont=\bfseries\sffamily,
    bodyfont=\normalfont,
  }
\newkeytheoremstyle{defbox}
  {
    tcolorbox={
		colbacktitle=Red!20,
		colframe=Red!60,
		borderline west={1mm}{0pt}{Red!70!white},
		before title=\large\texttwemoji{clipboard}\normalsize\hspace{2mm},
    },
    notebraces={(}{)},
    headpunct={.},
    noteseparator={ },
    headfont=\bfseries\sffamily,
    bodyfont=\normalfont,
  }
\newkeytheoremstyle{examplebox}
  {
    tcolorbox={
		colbacktitle=Green!25,
		colframe=Green!70,
		borderline west={1mm}{0pt}{Green},
		before title=\large\texttwemoji{chart increasing}\normalsize\hspace{2mm},
    },
    notebraces={(}{)},
    headpunct={.},
    noteseparator={ },
    headfont=\bfseries\sffamily,
    bodyfont=\normalfont,
  }

\newkeytheoremstyle{remarkbox}
  {
    tcolorbox={
		colbacktitle=Black!15,
		colframe=Black!30,
		borderline west={1mm}{0pt}{Black!50},
		before title=\large\texttwemoji{information}\normalsize\hspace{2mm},
    },
    notebraces={(}{)},
    headpunct={.},
    noteseparator={ },
    headfont=\bfseries\sffamily,
    bodyfont=\normalfont,
  }
\newkeytheoremstyle{warningbox}
  {
    tcolorbox={
		colbacktitle=Goldenrod!60,
		colframe=Goldenrod!80!Black,
		borderline west={1mm}{0pt}{Goldenrod!70!Black},
		before title=\large\texttwemoji{warning}\normalsize\hspace{2mm},
    },
    notebraces={(}{)},
    headpunct={.},
    noteseparator={ },
    headfont=\bfseries\sffamily,
    bodyfont=\normalfont,
  }

\newkeytheoremstyle{exercisebox}
  {
    tcolorbox={
    		detach title,
    		before upper={\tcbtitle\ \ },
		colback=White!93!Black,
		colframe=White!50!Black,
		leftrule = 1pt,
		rightrule = 1pt,
		toprule = 1pt,
		bottomrule = 1pt,
		boxsep = 1mm,
		before title=\large\texttwemoji{writing hand}\normalsize\hspace{2mm},
    },
    notebraces={(}{)},
    headpunct={.},
    noteseparator={ },
    headfont=\bfseries\sffamily,
    bodyfont=\normalfont,
  } 
  
  

\newkeytheorem{theorem}[name=Theorem, style=thmbox, parent=section]
\newkeytheorem{lemma}[name=Lemma, style=thmbox, sibling=theorem]
\newkeytheorem{proposition}[name=Proposition, style=thmbox, sibling=theorem]
\newkeytheorem{corollary}[name=Corollary, style=thmbox, sibling=theorem]
\newkeytheorem{fact}[name=Fact, style=thmbox, sibling=theorem]


\newkeytheorem{definition}[name = Definition, style=defbox, sibling=theorem]
\newkeytheorem{setup}[name = Setup, style=defbox, sibling=theorem]
\newkeytheorem{axiom}[name = Axiom, style=defbox, sibling=theorem]

\newkeytheorem{example}[name=Example, style=examplebox, sibling=theorem]
\newkeytheorem{conjecture}[name=Conjecture, style=examplebox, sibling=theorem]
\newkeytheorem{idea}[name=Idea, style=examplebox, sibling=theorem]

\newkeytheorem{remark}[name=Remark, style=remarkbox, sibling=theorem]
\newkeytheorem{warn}[name = Warning, style=warningbox, sibling=theorem]
\newkeytheorem{abuse}[name = Abuse of Notation, style=warningbox, sibling=theorem]

\newkeytheorem{exercise}[name = Exercise, style=exercisebox, sibling=theorem]


\newkeytheorem{lemma*}[name=Lemma, style=thmbox, numbered=no]
\newkeytheorem{proposition*}[name=Proposition, style=thmbox, numbered=no]
\newkeytheorem{corollary*}[name=Corollary, style=thmbox, numbered=no]
\newkeytheorem{fact*}[name=Fact, style=thmbox, numbered=no]

\newkeytheorem{definition*}[name = Definition, style=defbox, sibling=theorem, numbered=no]
\newkeytheorem{setup*}[name = Setup, style=defbox, sibling=theorem, numbered=no]

\newkeytheorem{example*}[name=Example, style=examplebox, sibling=theorem, numbered=no]
\newkeytheorem{idea*}[name=Idea, style=examplebox, sibling=theorem, numbered=no]

\newkeytheorem{remark*}[name=Remark, style=remarkbox, sibling=theorem, numbered=no]
\newkeytheorem{warn*}[name = Warning, style=defbox, sibling=theorem, numbered=no]

\theoremstyle{definition}
\newtheorem{problem}{Problem}[chapter]
\renewcommand\theproblem{\thechapter.\Alph{problem}}



